\documentclass[11pt,a4paper]{article}
\usepackage[T1]{fontenc}
\usepackage[utf8]{inputenc}
\usepackage{amsmath,amssymb,xcolor}
\usepackage[margin=25mm]{geometry}
\usepackage[colorlinks=true,urlcolor=blue]{hyperref}
\definecolor{nexorablue}{HTML}{2459B5}
\setlength{\parindent}{0pt}
\setlength{\parskip}{7pt}
\setlength{\emergencystretch}{2em}
\renewcommand{\familydefault}{\sfdefault}
\begin{document}

{\small\bfseries\color{nexorablue} NEXORA CPGE}\par
{\small MPSI / NEX-MES-02 / Corrigé}\par
{\LARGE\bfseries Résolution et étalonnage indépendants}\par
\vspace{0.25em}\hrule\vspace{0.6em}

\section*{Incertitudes élémentaires}
Pour l’arrondi, $u_r=0{,}050/\sqrt3=0{,}0289\,\mathrm V$. Pour l’étalonnage, $u_e=0{,}080/\sqrt3=0{,}0462\,\mathrm V$.

\section*{Budget}
L’indépendance donne $u=\sqrt{u_r^2+u_e^2}\simeq0{,}0545\,\mathrm V$. En adoptant deux chiffres significatifs pour $u$, on peut écrire $(4{,}600\pm0{,}054)\,\mathrm V$. L’étalonnage domine le budget.


\end{document}
